% =====================================================================
\chapter{Jerarquía de memoria y modelo roofline}
\label{tema:2}
% =====================================================================

Calcular es barato; traer los datos es caro. Este tema describe la jerarquía de memoria de un acelerador, define la intensidad aritmética de una operación y demuestra la cota del modelo \emph{roofline}, que decide si una operación está limitada por el cómputo o por la memoria. Después aplica el modelo al producto de matrices, al producto de matriz por vector y a la generación de texto, y analiza el \emph{tiling} como técnica para aumentar la intensidad. Se siguen el artículo de Williams, Waterman y Patterson~\cite{roofline}, el capítulo~4 de Hennessy y Patterson~\cite{hp}, el capítulo «All About Rooflines» de \emph{How to Scale Your Model}~\cite{scaling}, la clase~6 de CS336~\cite{cs336} y el resultado de Hong y Kung~\cite{hongkung}.

\section{La jerarquía de memoria}

Cada nivel de memoria de un acelerador es más grande y más lento que el anterior. La \cref{tab:jerarquia} resume los niveles de la A100; las cifras son aproximadas y proceden de la hoja de especificaciones y de la guía de ajuste de NVIDIA~\cite{a100}.

\begin{table}[h]
\centering\small
\begin{tabularx}{\textwidth}{@{}>{\raggedright\arraybackslash}p{3.3cm} >{\raggedright\arraybackslash}p{3.3cm} >{\raggedright\arraybackslash}p{2.6cm} >{\raggedright\arraybackslash}X@{}}
\toprule
Nivel & Capacidad & Ancho de banda & Uso \\
\midrule
Registros & 256~KB por SM & máximo & operandos de cada instrucción \\
SRAM (L1 y memoria compartida) & 192~KB por SM ($\approx 20$~MB en total) & decenas de TB/s agregados & bloques de datos gestionados por el programa (\emph{tiling}) \\
Caché L2 & 40~MB & intermedio & compartida por todos los SM \\
HBM & 80~GB & $\approx 2$~TB/s & pesos, activaciones y estado del optimizador \\
NVLink (entre GPU de un nodo) & --- & 600~GB/s & comunicación del paralelismo \\
Red entre nodos (InfiniBand) & --- & decenas de GB/s por GPU & el enlace más lento \\
\bottomrule
\end{tabularx}
\caption{Jerarquía de memoria y de comunicación de una A100 SXM de 80~GB (valores aproximados).}
\label{tab:jerarquia}
\end{table}

La HBM (\emph{High Bandwidth Memory}) es una pila de chips de DRAM montada junto al procesador: grande, pero lejana. La SRAM está dentro de cada SM: diminuta, pero muy rápida. Un \emph{kernel} eficiente carga un bloque de datos de la HBM a la SRAM, realiza con él todas las operaciones posibles y solo entonces escribe el resultado.

\section{Intensidad aritmética y modelo roofline}

\begin{definicion}[Intensidad aritmética]
Sea un \emph{kernel} que realiza $F$ FLOP y transfiere $Q$ bytes entre la memoria principal (HBM) y el procesador. Su \emph{intensidad aritmética} es $I=F/Q$, en FLOP por byte.
\end{definicion}

\begin{definicion}[Modelo roofline]
Sea un dispositivo con rendimiento máximo $\Ppico$ (FLOP/s) y ancho de banda de memoria $\Bmem$ (bytes/s). Su \emph{intensidad crítica} es $\Icrit=\Ppico/\Bmem$. El \emph{modelo roofline}~\cite{roofline} asigna a un \emph{kernel} de intensidad $I$ el rendimiento alcanzable
\[
P(I)=\min\bigl(\Ppico,\;I\,\Bmem\bigr).
\]
Se dice que el \emph{kernel} está \emph{limitado por la memoria} si $I<\Icrit$ y \emph{limitado por el cómputo} si $I>\Icrit$.
\end{definicion}

\begin{teorema}[Cota del roofline]\label{teo:roofline}
Sea un \emph{kernel} que realiza $F$ FLOP y transfiere $Q$ bytes en un dispositivo con parámetros $\Ppico$ y $\Bmem$. Su tiempo de ejecución cumple
\[
t\ \ge\ \max\Bigl(\frac{F}{\Ppico},\ \frac{Q}{\Bmem}\Bigr),
\qquad\text{y por tanto}\qquad
\frac{F}{t}\ \le\ \min\bigl(\Ppico,\ I\,\Bmem\bigr),\quad I=\frac{F}{Q}.
\]
\end{teorema}

La cota se alcanza solo si el cálculo y la transferencia se solapan por completo; si se ejecutasen en serie, el tiempo sería la suma $F/\Ppico+Q/\Bmem$, que es como mucho el doble del máximo. En ejes logarítmicos, $P(I)$ es una rampa de pendiente~1 seguida de un techo horizontal: de ahí el nombre (\cref{fig:roofline}).

\addtocounter{figure}{1}%

\addtocounter{definicion}{1}%

\section{Producto de matrices y producto de matriz por vector}

\begin{proposicion}[Intensidad del producto de matrices]\label{prop:intensidad-matmul}
Sea $C=AB$ con $A\in\R^{m\times k}$, $B\in\R^{k\times n}$, y sean $s$ los bytes por número. Si cada matriz se lee o escribe una sola vez, la intensidad es
\[
I=\frac{2mkn}{s\,(mk+kn+mn)}.
\]
En particular, para matrices cuadradas ($m=k=n$), $I=2n/(3s)$, que en bf16 ($s=2$) vale $n/3$.
\end{proposicion}

\begin{corolario}[Producto de matriz por vector]\label{cor:matvec}
Si $n=1$ (un solo vector) y $m,k\gg 1$, la intensidad es $I\approx 2/s$, que en bf16 vale $1$~FLOP/byte. Con $b$ vectores ($n=b\ll m,k$), $I\approx 2b/s$, que en bf16 vale $b$.
\end{corolario}

En el producto de matriz por vector cada peso leído se usa en una sola multiplicación. En la A100, $P\le 1\cdot 2=2$~TFLOP/s, un $0{,}6\,\%$ del máximo. Por eso la generación de texto token a token está limitada por la memoria (módulo~19): con lote~1, cada token exige un producto de matriz por vector con cada matriz de pesos.

\addtocounter{definicion}{1}%

\section{Tiling y cota inferior de tráfico}

Un producto de matrices grande no cabe en la SRAM. El tráfico efectivo con la HBM depende de cómo se recorra.

\begin{proposicion}[Tráfico del producto por bloques]\label{prop:tiling}
Sean $A,B\in\R^{n\times n}$ y sea $b$ un divisor de $n$. El algoritmo que, para cada bloque $b\times b$ de $C$, recorre los $n/b$ pares de bloques $b\times b$ de $A$ y de $B$ que contribuyen a él, los carga en la SRAM y acumula su producto, transfiere
\[
Q(b)=s\Bigl(\frac{2n^{3}}{b}+n^{2}\Bigr)\ \text{bytes},
\]
necesita una SRAM con capacidad para $3b^{2}$ números y tiene intensidad $I(b)=2n^{3}/Q(b)\approx b/s$ si $n\gg b$.
\end{proposicion}

La intensidad crece linealmente con el tamaño del bloque. Si la SRAM admite $M$ números, la condición $3b^{2}\le M$ da $b\le\sqrt{M/3}$ y un tráfico del orden de $n^{3}/\sqrt{M}$. Hong y Kung demostraron que este orden es óptimo:

\begin{teorema}[Hong y Kung, 1981]\label{teo:hongkung}
Todo algoritmo que calcula el producto de dos matrices $n\times n$ por el método clásico de $n^{3}$ productos con una memoria rápida de $M$ números transfiere $\Omega\bigl(n^{3}/\sqrt{M}\bigr)$ números entre la memoria rápida y la lenta.
\end{teorema}

La demostración, mediante el juego de los guijarros rojos y azules, puede consultarse en~\cite{hongkung}. El \cref{eje:2-tiling} pide aplicar la \cref{prop:tiling} a la SRAM de un SM.

\begin{observacion}
La \cref{prop:tiling} describe un solo nivel de la jerarquía. En una GPU hay varios: registros, SRAM y caché L2. La caché L2 (40~MB en la A100) permite que bloques de hilos vecinos reutilicen los mismos bloques de $A$ y $B$, de modo que el tráfico real con la HBM es menor que el que da la proposición con el $b$ de la SRAM. El principio es el mismo en cada nivel: cuanto mayor es el bloque que cabe, mayor es la reutilización.
\end{observacion}

\section{Utilización efectiva}

\begin{definicion}[Utilización de FLOP del modelo]\label{def:mfu}
Sea un entrenamiento que procesa $\tau$ tokens por segundo con un modelo de $N$ parámetros en $n_{\mathrm{d}}$ dispositivos de rendimiento máximo $\Ppico$. Su \emph{utilización de FLOP del modelo} es~\cite[§4.1 y apéndice~B]{palm}
\[
\MFU=\frac{6N\tau}{n_{\mathrm{d}}\,\Ppico}.
\]
\end{definicion}

El numerador cuenta solo los FLOP «útiles» del \cref{cor:6nd}; la recomputación de activaciones (\cref{tema:5}) no suma. Los entrenamientos grandes bien optimizados se sitúan típicamente entre el 30~\% y el 60~\% de MFU (PaLM 540B declaró un 46{,}2~\%~\cite[§4.1]{palm}); el resto se pierde en comunicación, en operaciones limitadas por la memoria y en esperas.

